%% GENERATED by tools/gen_latex.py -- edit content/agents08.py instead.
%% Build:  latexmk -pdf agents08.tex     (or: pdflatex agents08.tex, twice)
\documentclass[aspectratio=169,11pt,t]{beamer}
\usepackage{itbpro}

\coursetitle{Sistem banyak agen}
\coursesubtitle{Kapan memecah pekerjaan jadi beberapa agen benar-benar menolong --- dan variabel yang menentukan harganya, yang hampir tidak pernah diukur.}
\coursekicker{Chapter 8}
\coursesource{Grootendorst \& Alammar, An Illustrated Guide to AI Agents --- bab 8}
\courseduration{3 jam (2 sesi)}
\coursepresenter{Hendri Karisma}{Instructor}
\coursebrand{AI for Professional \textperiodcentered\ ITB \textperiodcentered\ Ch 8}

\title{Sistem banyak agen}
\author{Hendri Karisma}
\date{}

\begin{document}
\covertitle

\framekicker{Learning outcomes}
\begin{frame}[shrink=25]{Session Objectives}
\leadpar{By the end of this session, participants can:}
\begin{isteps}
\item \textbf{Menerapkan tiga pertanyaan} yang memutuskan sebuah pekerjaan pantas dipecah jadi beberapa agen.
\item \textbf{Menghitung biaya memecah}, dan menyebutkan variabel yang sebenarnya menentukannya.
\item \textbf{Menyebutkan tiga pola orkestrasi} dan bentuk kegagalan khas masing-masing.
\item \textbf{Menjelaskan apa yang hilang di tiap serah-terima}, dan cara membatasinya.
\item \textbf{Membedakan protokol ke dalam dari protokol menyilang}, dan pertanyaan kepercayaan yang muncul pada yang kedua.
\item \textbf{Menelusuri kegagalan pada sistem banyak agen} dengan jejak per agen.
\end{isteps}

\vskip 6pt
\vskip 4pt
\reslink{Site}{Course home}{https://hendrikarisma.my.id/rs/ai-products-course/index.html}
\reslink{Repo}{ai-agentic-demo --- empat kasus multi-agent, dengan tabel biayanya}{https://github.com/situkangsayur/ai-agentic-demo}
\reslink{Book}{Maarten Grootendorst \& Jay Alammar, An Illustrated Guide to AI Agents (O'Reilly Media, Early Release)}{https://www.oreilly.com/library/view/an-illustrated-guide/9798341662681/}
\vskip 2pt
\end{frame}

\coursesection{01}{Kapan memecah}{Tiga pertanyaan, dan bagan organisasi bukan salah satunya.}

\framekicker{Memutuskan}
\begin{frame}[shrink=25]{Tiga bentuk yang membuat pemecahan masuk akal}
\begin{minipage}[t]{0.3253\textwidth}
\begin{icardgood}
\cardh{Benar-benar sejajar}
Sub-tugas yang tidak saling membutuhkan hasil. Menilai lima pemasok, meringkas dua puluh dokumen. Ini satu-satunya pemecahan yang tidak perlu dibela.
\end{icardgood}
\end{minipage}%
\hfill
\begin{minipage}[t]{0.3253\textwidth}
\begin{icardgood}
\cardh{Batas izin}
Satu agen boleh membaca data sensitif, satu lagi boleh menulis, dan tak satu pun boleh keduanya. Pemisahannya \textbf{keamanan}, bukan efisiensi.
\end{icardgood}
\end{minipage}%
\hfill
\begin{minipage}[t]{0.3253\textwidth}
\begin{icardgood}
\cardh{Kumpulan alat yang tak berhubungan}
Kalau satu daftar alat sudah terlalu besar dan terbelah rapi jadi dua yang tak saling pakai.
\end{icardgood}
\end{minipage}%
{\fontsize{9.4}{12.6}\selectfont\color{fg2}Yang \textbf{tidak} ada di daftar ini: meniru struktur tim manusia. \textbackslash\{\}u201cPeneliti, penulis, kritikus\textbackslash\{\}u201d rapi di papan tulis dan tidak menjawab satu pun dari tiga pertanyaan di atas.\par}\vskip 4pt
\end{frame}

\framekicker{Memutuskan}
\begin{frame}[shrink=25]{Uji yang paling cepat: apakah mereka butuh konteks yang sama?}
{\fontsize{9.4}{12.6}\selectfont\color{fg2}Kalau dua \textbackslash\{\}u201cagen\textbackslash\{\}u201d membutuhkan sebagian besar konteks yang sama untuk bekerja, memecahnya berarti \textbf{mengirim konteks itu dua kali} dan kehilangan sebagian di antaranya.\par}\vskip 4pt
\begin{minipage}[t]{0.485\textwidth}
{\fontsize{9.4}{12.6}\selectfont\color{fg2}\textbf{Tanda pemecahan yang salah}\par}\vskip 4pt
\begin{ibullets}
\item Serah-terimanya hampir sepanjang percakapannya
\item Agen kedua sering meminta hal yang sudah diketahui agen pertama
\item Keduanya memakai daftar alat yang sama
\end{ibullets}

\end{minipage}%
\hfill
\begin{minipage}[t]{0.485\textwidth}
{\fontsize{9.4}{12.6}\selectfont\color{fg2}\textbf{Tanda pemecahan yang benar}\par}\vskip 4pt
\begin{ibullets}
\item Serah-terimanya pendek dan bentuknya tetap
\item Daftar alatnya berbeda
\item Salah satunya bisa berjalan tanpa yang lain
\end{ibullets}

\end{minipage}%
\begin{iband}[signal]
Uji satu kalimat: \textbf{kalau serah-terimanya harus panjang supaya agen berikutnya bisa bekerja, keduanya sebenarnya satu agen} yang sedang dipaksa berbicara lewat surat.
\end{iband}
\end{frame}

\coursesection{02}{Harganya, dihitung}{Dan jawabannya bukan yang saya duga sebelum menghitungnya.}

\framekicker{Biaya}
\begin{frame}[shrink=25]{Dugaan yang wajar, dan hitungan yang membantahnya}
\begin{center}
\includegraphics[width=0.97\linewidth,height=0.56\textheight,keepaspectratio]{figs/agents08-splitcost.tex.pdf}
\end{center}
\vskip 3pt{\color{ink3}\fontsize{7.4}{9.4}\selectfont Satu agen lawan tiga agen yang membaginya, dihitung. Langkahi menurut ukuran serah-terima.\par}
{\fontsize{9.4}{12.6}\selectfont\color{fg2}\textbackslash\{\}u201cMemecah pasti lebih mahal, karena ada biaya koordinasi\textbackslash\{\}u201d --- masuk akal, dan \textbf{salah} untuk serah-terima yang kecil. Dengan serah-terima 600 token, tiga agen justru lebih murah di semua panjang tugas.\par}\vskip 4pt
\end{frame}

\framekicker{Biaya}
\begin{frame}[shrink=25]{Variabel yang menentukan, dan yang tidak diukur siapa pun}
{\fontsize{9.4}{12.6}\selectfont\color{fg2}Yang menentukan bukan jumlah agennya, melainkan \textbf{ukuran serah-terimanya}. Serah-terima kecil memecah kuadrat dan menghemat; serah-terima besar dibawa ke setiap agen sesudahnya dan membengkak.\par}\vskip 4pt
\begin{itable}{>{\raggedright\arraybackslash}p{0.2256\textwidth}>{\raggedright\arraybackslash}p{0.188\textwidth}>{\raggedright\arraybackslash}p{0.188\textwidth}>{\raggedright\arraybackslash}p{0.3384\textwidth}}
\thead{Serah-terima} & \thead{6 langkah} & \thead{12 langkah} & \thead{Artinya} \\ \midrule
600 token & 0,93$\times$ & 0,81$\times$ & Memecah menghemat, makin panjang makin hemat \\
1 500 token & 1,07$\times$ & 0,92$\times$ & Bergantung panjang tugasnya \\
3 000 token & 1,31$\times$ & 1,11$\times$ & Memecah mahal --- dan inilah yang biasanya terjadi \\
\end{itable}
\begin{iband}[signal]
Baris ketiga yang paling sering nyata, sebab serah-terima cenderung membengkak: setiap kali agen berikutnya kekurangan informasi, orang menambahkan sesuatu ke ringkasannya --- \hilite{dan tidak pernah mengurangi}
\end{iband}
\end{frame}

\framekicker{Biaya}
\begin{frame}[shrink=25]{Yang benar-benar mahal bukan tokennya}
\begin{minipage}[t]{0.494\textwidth}
\begin{icardbad}
\cardh{Informasi yang hilang di serah-terima}
Ini kerugian yang sesungguhnya, dan ia tidak muncul di tagihan mana pun. Bab 4 sudah menghitung bentuknya: peringkasan kehilangan angka dan pengecualian lebih dulu.
\end{icardbad}
\end{minipage}%
\hfill
\begin{minipage}[t]{0.494\textwidth}
\begin{icardbad}
\cardh{Sulit ditelusuri}
Satu jejak jadi tiga jejak, dan pertanyaan \textbackslash\{\}u201cdi mana salahnya\textbackslash\{\}u201d butuh membaca ketiganya beserta sambungannya.
\end{icardbad}
\end{minipage}%


\vskip 4pt

\begin{minipage}[t]{0.494\textwidth}
\begin{icard}
\cardh{Waktu dinding bertambah}
Kecuali kalau agennya benar-benar berjalan sejajar --- dan sebagian besar pemecahan tidak.
\end{icard}
\end{minipage}%
\hfill
\begin{minipage}[t]{0.494\textwidth}
\begin{icard}
\cardh{Lebih banyak yang bisa gagal}
Tiga gelung, tiga anggaran, tiga cara berhenti di tempat yang salah.
\end{icard}
\end{minipage}%
{\fontsize{9.4}{12.6}\selectfont\color{fg2}Karena itu argumen melawan pemecahan yang berlebihan tetap berlaku --- hanya saja \textbf{alasannya bukan biaya token}, dan menyebut alasan yang salah membuat perdebatannya tidak selesai.\par}\vskip 4pt
\end{frame}

\coursesection{03}{Pola orkestrasi}{Tiga bentuk, tiga cara gagal.}

\framekicker{Pola}
\begin{frame}[shrink=25]{Tiga bentuk yang mencakup hampir semua sistem nyata}
\begin{center}
\includegraphics[width=0.92\linewidth,height=0.60\textheight,keepaspectratio]{figs/agents08-patterns.tex.pdf}
\end{center}
\vskip 3pt{\color{ink3}\fontsize{7.4}{9.4}\selectfont Penyelia, berantai, dan kelompok. Yang ketiga paling menarik dibaca dan paling jarang pantas dipakai di produksi: tanpa pemilik keputusan, tidak ada yang bertanggung jawab menghentikannya.\par}
\begin{itable}{>{\raggedright\arraybackslash}p{0.188\textwidth}>{\raggedright\arraybackslash}p{0.376\textwidth}>{\raggedright\arraybackslash}p{0.376\textwidth}}
\thead{Pola} & \thead{Cocok untuk} & \thead{Kegagalan khasnya} \\ \midrule
Penyelia & Sub-tugas sejajar dengan satu pemilik hasil & Penyelia jadi leher botol dan konteksnya membengkak \\
Berantai & Tahapan yang urutannya tetap & Kesalahan di tahap awal dibangun oleh semua tahap sesudahnya \\
Kelompok & Diskusi terbuka tanpa pemilik jelas & Tidak ada yang menghentikannya; biaya sulit diramalkan \\
\end{itable}
\end{frame}

\framekicker{Pola}
\begin{frame}[shrink=25]{Penyelia adalah kode, bukan agen --- kalau bisa}
{\fontsize{9.4}{12.6}\selectfont\color{fg2}Pertanyaan yang jarang ditanyakan: apakah penyelianya perlu berupa model sama sekali? Kalau urutan dan syaratnya sudah diketahui, penyelia terbaik adalah \textbf{fungsi biasa}.\par}\vskip 4pt
\begin{minipage}[t]{0.485\textwidth}
{\fontsize{9.4}{12.6}\selectfont\color{fg2}\textbf{Penyelia berupa kode}\par}\vskip 4pt
\begin{ibullets}
\item Deterministik dan gratis
\item Bisa diuji seperti kode biasa
\item Tidak bisa dibujuk
\end{ibullets}

\end{minipage}%
\hfill
\begin{minipage}[t]{0.485\textwidth}
{\fontsize{9.4}{12.6}\selectfont\color{fg2}\textbf{Penyelia berupa model}\par}\vskip 4pt
\begin{ibullets}
\item Perlu kalau pembagian kerjanya benar-benar bervariasi
\item Menambah satu gelung, satu anggaran, satu tempat gagal
\end{ibullets}

\end{minipage}%
\begin{iband}[signal]
Ini bentuk lain dari uji di Bab 1: \textbf{kalau langkahnya sudah diketahui, yang Anda butuhkan alur tetap} --- dan itu berlaku juga untuk lapisan koordinasinya.
\end{iband}
\end{frame}

\coursesection{04}{Serah-terima}{Satu-satunya jalan informasi berpindah, dan tempat sebagian besar hilang.}

\framekicker{Serah-terima}
\begin{frame}[shrink=25]{Agen berikutnya hanya melihat apa yang Anda kirimkan}
\begin{center}
\includegraphics[width=0.92\linewidth,height=0.60\textheight,keepaspectratio]{figs/agents08-handoff.tex.pdf}
\end{center}
\vskip 3pt{\color{ink3}\fontsize{7.4}{9.4}\selectfont Konteks penuh tinggal di agen pertama; agen kedua melihat ringkasannya.\par}
{\fontsize{9.4}{12.6}\selectfont\color{fg2}Ini pengulangan Bab 4 dalam bentuk yang lebih tajam: apa pun yang tidak masuk ke serah-terima \textbf{tidak ada} bagi agen berikutnya --- dan tidak ada mekanisme yang memberitahunya bahwa sesuatu hilang.\par}\vskip 4pt
\begin{iband}[signal]
Yang paling sering hilang tetap yang sama: \textbf{angka, pengecualian, dan hal yang belum diputuskan.} Ketiganya justru yang paling menentukan keputusan berikutnya.
\end{iband}
\end{frame}

\framekicker{Serah-terima}
\begin{frame}[shrink=25]{Bentuk serah-terima yang tidak membengkak}
\begin{center}
\includegraphics[width=0.92\linewidth,height=0.60\textheight,keepaspectratio]{figs/agents08-hshape.tex.pdf}
\end{center}
\vskip 3pt{\color{ink3}\fontsize{7.4}{9.4}\selectfont Empat medan yang tetap --- bukan prosa bebas.\par}
{\fontsize{9.4}{12.6}\selectfont\color{fg2}Prosa bebas selalu membengkak, sebab tiap kali agen berikutnya kekurangan sesuatu, orang menambahkan satu kalimat lagi --- dan tidak pernah mengurangi. Medan yang tetap membuat pertumbuhan itu \textbf{terlihat}.\par}\vskip 4pt
\end{frame}

\framekicker{Serah-terima}
\begin{frame}[shrink=25]{\ldots{}dan empat aturan yang menjaganya tetap kecil}
\begin{isteps}
\item \{'h': 'Terstruktur, bukan prosa', 'p': 'Medan yang tetap: temuan, angka beserta sumbernya, yang masih terbuka, dan apa yang diminta dari agen berikutnya.'\}
\item \{'h': 'Rujukan, bukan salinan', 'p': 'Hasil alat yang besar tinggal di tempatnya; kirimkan pengenalnya, dan beri agen berikutnya alat untuk membacanya.'\}
\item \{'h': 'Batasi ukurannya, dan pantau', 'p': 'Kalau serah-terima tumbuh dari rilis ke rilis, pemecahannya sedang menuju baris ketiga tabel biaya tadi.'\}
\item \{'h': 'Tulis apa yang TIDAK diketahui', 'p': 'Ketidakpastian yang tidak diteruskan berubah jadi kepastian palsu: agen kedua yang tidak diberi tahu bahwa sebuah angka adalah taksiran akan memperlakukannya sebagai fakta.'\}
\end{isteps}
\end{frame}

\framekicker{Peran}
\begin{frame}[shrink=25]{Memberi agen ``peran'' lebih jarang menolong daripada yang dikira}
{\fontsize{9.4}{12.6}\selectfont\color{fg2}Pola yang paling banyak disalin: beri tiap agen persona --- peneliti, penulis, kritikus --- dengan harapan spesialisasi menghasilkan mutu.\par}\vskip 4pt
{\fontsize{9.4}{12.6}\selectfont\color{fg2}Yang sebenarnya membuat agen berperilaku berbeda bukan kalimat perannya, melainkan \textbf{alat yang dimilikinya dan konteks yang dilihatnya}. Dua agen dengan alat dan konteks identik akan berperilaku hampir sama, betapapun berbeda kalimat perannya.\par}\vskip 4pt
\begin{iband}[signal]
Uji yang menyelesaikan perdebatan ini dengan cepat: \textbf{tukar kalimat peran di antara dua agen dan jalankan lagi.} Kalau hasilnya nyaris sama, perannya tidak sedang mengerjakan apa pun.
\end{iband}
\end{frame}

\framekicker{Peran}
\begin{frame}[shrink=25]{Spesialisasi yang benar-benar bekerja}
\begin{minipage}[t]{0.3253\textwidth}
\begin{icardgood}
\cardh{Alat yang berbeda}
Agen yang hanya punya alat baca benar-benar tidak bisa menulis. Ini spesialisasi yang berupa struktur.
\end{icardgood}
\end{minipage}%
\hfill
\begin{minipage}[t]{0.3253\textwidth}
\begin{icardgood}
\cardh{Konteks yang berbeda}
Agen yang tidak melihat data pribadi tidak bisa membocorkannya. Juga struktur.
\end{icardgood}
\end{minipage}%
\hfill
\begin{minipage}[t]{0.3253\textwidth}
\begin{icardbad}
\cardh{Kalimat peran saja}
Mempengaruhi gaya, hampir tidak mempengaruhi keputusan --- dan bisa dibujuk.
\end{icardbad}
\end{minipage}%
{\fontsize{9.4}{12.6}\selectfont\color{fg2}Pola yang sama muncul di tiap bab modul ini: \textbf{yang menentukan perilaku adalah apa yang mungkin dilakukan, bukan apa yang diminta.} Peran adalah permintaan; alat adalah kemungkinan.\par}\vskip 4pt
\end{frame}

\framekicker{Komunikasi}
\begin{frame}[shrink=25]{Sinkron atau tidak, dan siapa yang menunggu siapa}
\begin{minipage}[t]{0.485\textwidth}
{\fontsize{9.4}{12.6}\selectfont\color{fg2}\textbf{Berurutan (sinkron)}\par}\vskip 4pt
{\fontsize{9.4}{12.6}\selectfont\color{fg2}Agen berikutnya menunggu yang sebelumnya. Mudah ditelusuri, mudah dibatasi, dan waktunya menumpuk.\par}\vskip 4pt
{\fontsize{9.4}{12.6}\selectfont\color{fg2}Ini yang benar untuk hampir semua alur kerja bisnis.\par}\vskip 4pt

\end{minipage}%
\hfill
\begin{minipage}[t]{0.485\textwidth}
{\fontsize{9.4}{12.6}\selectfont\color{fg2}\textbf{Sejajar (asinkron)}\par}\vskip 4pt
{\fontsize{9.4}{12.6}\selectfont\color{fg2}Beberapa berjalan bersamaan lalu hasilnya disatukan. Menghemat waktu dinding, dan menciptakan pertanyaan urutan.\par}\vskip 4pt
{\fontsize{9.4}{12.6}\selectfont\color{fg2}Hanya aman kalau semuanya \textbf{alat baca}.\par}\vskip 4pt

\end{minipage}%
\begin{iband}[signal]
Aturan yang tidak bisa ditawar: \textbf{jangan menjalankan dua alat tulis secara sejajar.} Kalau keduanya menyentuh keadaan yang sama, hasilnya bergantung pada urutan yang tidak Anda kendalikan.
\end{iband}
\end{frame}

\framekicker{Komunikasi}
\begin{frame}[shrink=25]{Keadaan bersama adalah tempat bug berkumpul}
{\fontsize{9.4}{12.6}\selectfont\color{fg2}Begitu beberapa agen menulis ke tempat yang sama --- antrean, papan catatan, ingatan bersama --- sistem Anda punya semua masalah yang dikenal sistem terdistribusi, ditambah satu komponen yang tidak deterministik.\par}\vskip 4pt
\begin{isteps}
\item \{'h': 'Satu penulis per bagian keadaan', 'p': 'Kalau dua agen boleh menulis hal yang sama, tentukan siapa yang menang --- di kode, sebelum kejadian.'\}
\item \{'h': 'Tulis dengan kunci idempotensi', 'p': 'Bab 5 sudah menyebutnya, dan di sini ia jadi wajib, bukan dianjurkan.'\}
\item \{'h': 'Beri urutan yang bisa dibaca', 'p': 'Stempel waktu dan pengenal proses di tiap tulisan, atau rekonstruksi kejadian jadi mustahil.'\}
\end{isteps}
\begin{iband}[signal]
Temuan yang tercatat di repo demo persis dari kelas ini: antrean yang semula sebuah dict berjalan sempurna sampai peladen alat jadi proses sendiri --- \textbf{kotaknya benar, panahnya benar, dan keadaannya ada di tempat yang salah.}
\end{iband}
\end{frame}

\coursesection{05}{Protokol antar agen}{Ke dalam dan menyilang adalah dua masalah yang berbeda.}

\framekicker{Protokol}
\begin{frame}[shrink=25]{Ke dalam sistem sendiri, dan menyilang ke milik orang lain}
\begin{center}
\includegraphics[width=0.92\linewidth,height=0.60\textheight,keepaspectratio]{figs/agents08-a2a.tex.pdf}
\end{center}
\vskip 3pt{\color{ink3}\fontsize{7.4}{9.4}\selectfont Dua arah, dua pertanyaan kepercayaan yang berbeda.\par}
{\fontsize{9.4}{12.6}\selectfont\color{fg2}Protokol alat (Bab 5) menjawab \textbackslash\{\}u201cbagaimana agen memanggil sistem saya\textbackslash\{\}u201d. Protokol antar agen menjawab \textbackslash\{\}u201cbagaimana agen saya bicara dengan agen \textbf{yang bukan milik saya}\textbackslash\{\}u201d --- dan itu pertanyaan yang berbeda sepenuhnya.\par}\vskip 4pt
\begin{iband}[signal]
Perbedaannya bukan teknis melainkan \textbf{kepercayaan}: di dalam sistem sendiri, Anda tahu apa yang ada di ujung sana. Menyilang, yang datang kembali adalah masukan tidak tepercaya dari pihak yang tidak Anda kendalikan.
\end{iband}
\end{frame}

\framekicker{Protokol}
\begin{frame}[shrink=25]{Pertanyaan yang harus dijawab sebelum agen bicara ke luar}
\begin{minipage}[t]{0.494\textwidth}
\begin{icardaccent}
\cardh{Atas nama siapa?}
Identitas pengguna akhir harus ikut, atau agen di seberang tidak bisa memeriksa izin apa pun.
\end{icardaccent}
\end{minipage}%
\hfill
\begin{minipage}[t]{0.494\textwidth}
\begin{icardaccent}
\cardh{Apa yang boleh keluar?}
Serah-terima ke agen pihak lain adalah \textbf{perpindahan data}, dengan semua kewajiban yang menyertainya.
\end{icardaccent}
\end{minipage}%


\vskip 4pt

\begin{minipage}[t]{0.494\textwidth}
\begin{icard}
\cardh{Siapa yang menghentikannya?}
Gelung yang melintasi dua organisasi punya dua anggaran yang tidak saling tahu.
\end{icard}
\end{minipage}%
\hfill
\begin{minipage}[t]{0.494\textwidth}
\begin{icard}
\cardh{Siapa yang menyimpan jejaknya?}
Kalau keputusan diambil bersama, buktinya harus bisa disatukan kembali.
\end{icard}
\end{minipage}%
{\fontsize{9.4}{12.6}\selectfont\color{fg2}Untuk bank dan lembaga yang diawasi, keempatnya \textbf{mendahului} pertanyaan teknis mana pun tentang format pesan --- dan jawaban ketiganya jarang ada di spesifikasi protokol.\par}\vskip 4pt
\end{frame}

\framekicker{Protokol}
\begin{frame}[shrink=25]{Menemukan agen lain, dan mempercayainya}
{\fontsize{9.4}{12.6}\selectfont\color{fg2}Begitu agen boleh bicara ke luar, muncul dua pertanyaan yang tidak ada di dalam sistem sendiri: \textbf{bagaimana ia tahu agen lain itu ada}, dan \textbf{kenapa ia percaya pada jawabannya.}\par}\vskip 4pt
\begin{minipage}[t]{0.494\textwidth}
\begin{icard}
\cardh{Menemukan}
Daftar kemampuan yang bisa dibaca mesin --- apa yang bisa dikerjakan agen itu, dengan masukan apa. Bentuknya mirip skema alat, dan alasannya sama.
\end{icard}
\end{minipage}%
\hfill
\begin{minipage}[t]{0.494\textwidth}
\begin{icardaccent}
\cardh{Mempercayai}
Jawaban dari agen pihak lain adalah \textbf{masukan tidak tepercaya}, setara hasil pengambilan dokumen. Ia divalidasi, bukan diteruskan.
\end{icardaccent}
\end{minipage}%
\begin{iband}[signal]
Kesalahan yang mudah dibuat dan mahal: memperlakukan jawaban agen mitra sebagai fakta karena ia datang lewat saluran resmi. \hilite{Saluran yang resmi tidak membuat isinya benar}
\end{iband}
\end{frame}

\framekicker{Protokol}
\begin{frame}[shrink=25]{Batas organisasi mengubah pertanyaannya jadi pertanyaan kontrak}
{\fontsize{9.4}{12.6}\selectfont\color{fg2}Di dalam satu perusahaan, ``agen A memanggil agen B'' adalah keputusan arsitektur. Melintasi perusahaan, ia keputusan \textbf{hukum} yang kebetulan diwujudkan dengan kode.\par}\vskip 4pt
\begin{itable}{>{\raggedright\arraybackslash}p{0.2632\textwidth}>{\raggedright\arraybackslash}p{0.3196\textwidth}>{\raggedright\arraybackslash}p{0.3572\textwidth}}
\thead{Pertanyaan} & \thead{Di dalam} & \thead{Menyilang} \\ \midrule
Data apa yang keluar? & Kebijakan internal & Perjanjian pemrosesan data, dan catatannya \\
Siapa bertanggung jawab? & Satu tim & Harus tertulis sebelum sambungan pertama \\
Kalau salah? & Perbaiki dan jalankan lagi & Siapa yang menanggung kerugiannya? \\
Jejaknya di mana? & Satu tempat & Dua tempat, dan harus bisa disatukan \\
\end{itable}
\begin{iband}[signal]
Untuk lembaga yang diawasi, keempat baris kanan ini \textbf{mendahului} pilihan protokol --- dan tidak satu pun dijawab oleh spesifikasi teknis mana pun.
\end{iband}
\end{frame}

\framekicker{Simulasi}
\begin{frame}[shrink=25]{Perilaku yang muncul, dan cara membacanya dengan benar}
{\fontsize{9.4}{12.6}\selectfont\color{fg2}Sekelompok agen dengan aturan sederhana bisa menghasilkan pola yang tidak diprogramkan siapa pun --- pembagian kerja, pengulangan, bahkan sesuatu yang terlihat seperti kesepakatan.\par}\vskip 4pt
\begin{minipage}[t]{0.485\textwidth}
{\fontsize{9.4}{12.6}\selectfont\color{fg2}\textbf{Kesimpulan yang sah}\par}\vskip 4pt
{\fontsize{9.4}{12.6}\selectfont\color{fg2}Aturan yang saya beri menghasilkan pola ini pada sistem ini. Itu temuan tentang \textbf{rancangan saya}.\par}\vskip 4pt

\end{minipage}%
\hfill
\begin{minipage}[t]{0.485\textwidth}
{\fontsize{9.4}{12.6}\selectfont\color{fg2}\textbf{Kesimpulan yang tidak sah}\par}\vskip 4pt
{\fontsize{9.4}{12.6}\selectfont\color{fg2}Pola ini akan muncul pada manusia, atau pada sistem lain, atau berarti agennya ``memahami'' sesuatu.\par}\vskip 4pt

\end{minipage}%
\begin{iband}[signal]
Perbedaan keduanya bukan kehati-hatian akademis. Sistem yang dibangun di atas kesimpulan kolom kanan akan \textbf{gagal justru pada kasus yang paling penting}, sebab dasarnya tidak pernah ada.
\end{iband}
\end{frame}

\coursesection{06}{Simulasi dan riset}{Arah yang menarik, dan garis yang memisahkannya dari produksi.}

\framekicker{Simulasi}
\begin{frame}[shrink=25]{Banyak agen sebagai alat penelitian, bukan sebagai produk}
{\fontsize{9.4}{12.6}\selectfont\color{fg2}Sebagian penggunaan paling menarik dari sistem banyak agen bukan menyelesaikan tugas, melainkan \textbf{mensimulasikan perilaku}: sekelompok agen dengan tujuan berbeda dijalankan bersama, lalu polanya diamati.\par}\vskip 4pt
\begin{minipage}[t]{0.485\textwidth}
{\fontsize{9.4}{12.6}\selectfont\color{fg2}\textbf{Yang berguna dari sana}\par}\vskip 4pt
\begin{ibullets}
\item Menguji rancangan pasar atau kebijakan sebelum diterapkan
\item Membangkitkan kasus uji yang tidak terpikir
\item Melihat perilaku yang muncul dari aturan sederhana
\end{ibullets}

\end{minipage}%
\hfill
\begin{minipage}[t]{0.485\textwidth}
{\fontsize{9.4}{12.6}\selectfont\color{fg2}\textbf{Yang tidak boleh disimpulkan}\par}\vskip 4pt
\begin{ibullets}
\item Bahwa manusia akan berperilaku begitu
\item Bahwa hasilnya bisa dipakai sebagai bukti
\item Bahwa \textbackslash\{\}u201cagen sepakat\textbackslash\{\}u201d berarti sesuatu itu benar
\end{ibullets}

\end{minipage}%
\begin{iband}[signal]
Kesepakatan antar agen yang dibangun dari model yang sama adalah \textbf{kesepakatan antara satu model dan dirinya sendiri}. Bab 3 sudah menyebut bentuk kekeliruan yang sama pada suara terbanyak.
\end{iband}
\end{frame}

\framekicker{Riset}
\begin{frame}[shrink=25]{Agen riset mendalam: bentuk yang paling banyak dipakai orang}
{\fontsize{9.4}{12.6}\selectfont\color{fg2}Satu penerapan banyak agen yang benar-benar berguna dan sudah dipakai luas: memecah satu pertanyaan riset jadi beberapa pencarian yang berjalan \textbf{sejajar}, lalu menyatukan hasilnya.\par}\vskip 4pt
{\fontsize{9.4}{12.6}\selectfont\color{fg2}Perhatikan bahwa ini persis bentuk pertama di daftar \textbackslash\{\}u201ckapan memecah\textbackslash\{\}u201d: sub-tugas yang benar-benar tidak saling membutuhkan. Ia berhasil karena alasan yang bisa disebut, bukan karena banyak agen terdengar hebat.\par}\vskip 4pt
\begin{iband}[signal]
Dan syarat mutunya sama seperti bab 4: \textbf{tiap temuan membawa sumbernya.} Ringkasan riset tanpa kutipan yang bisa diperiksa adalah prosa yang meyakinkan, bukan riset.
\end{iband}
\end{frame}

\framekicker{Riset}
\begin{frame}[shrink=25]{Model teori pikiran, dan kenapa hati-hati membacanya}
{\fontsize{9.4}{12.6}\selectfont\color{fg2}Sebagian penelitian menguji apakah agen bisa memperhitungkan \textbf{apa yang diketahui agen lain} --- informasi yang dimiliki pihak lain berbeda dari yang dimilikinya sendiri.\par}\vskip 4pt
{\fontsize{9.4}{12.6}\selectfont\color{fg2}Ini berguna secara teknis: agen yang menyusun serah-terima harus memperkirakan apa yang belum diketahui penerimanya. Itu masalah rekayasa yang nyata, dan bisa diperbaiki dengan bentuk pesan yang tetap.\par}\vskip 4pt
\begin{iband}[signal]
Yang perlu hati-hati adalah bahasanya. Sistem yang lulus uji semacam itu \textbf{memodelkan informasi}, dan menyebutnya ``memahami orang lain'' menggeser harapan pemakainya ke tempat yang tidak didukung apa pun.
\end{iband}
\end{frame}

\framekicker{Kerangka kerja}
\begin{frame}[shrink=25]{Memilih kerangka kerja, dan pertanyaan yang lebih penting}
{\fontsize{9.4}{12.6}\selectfont\color{fg2}Ada banyak kerangka kerja untuk merangkai beberapa agen, dan pilihannya jauh kurang menentukan daripada yang diperdebatkan orang.\par}\vskip 4pt
\begin{itable}{>{\raggedright\arraybackslash}p{0.282\textwidth}>{\raggedright\arraybackslash}p{0.3196\textwidth}>{\raggedright\arraybackslash}p{0.3384\textwidth}}
\thead{Yang biasanya ditanyakan} & \thead{Yang lebih menentukan} & \thead{Kenapa} \\ \midrule
Kerangka kerja mana? & Bentuk serah-terimanya & Menentukan biaya dan mutu; kerangka kerja tidak \\
Bagaimana polanya? & Apakah memang perlu dipecah & Pola terbaik untuk pemecahan yang salah tetap salah \\
Bisa asinkron? & Apakah ada alat tulis di dalamnya & Sejajar aman untuk baca, berbahaya untuk tulis \\
Ada UI-nya? & Apakah jejaknya per agen & Tanpa itu, tidak ada yang bisa ditelusuri \\
\end{itable}
\begin{iband}[signal]
Dan satu sifat yang pantas diperiksa pada kerangka kerja mana pun: \textbf{bisakah Anda melihat konteks persis yang dikirim ke model di tiap giliran?} Kalau tidak, Anda tidak bisa memperbaiki hal yang paling sering perlu diperbaiki.
\end{iband}
\end{frame}

\framekicker{Menguji}
\begin{frame}[shrink=25]{Menguji sistem banyak agen: uji sambungannya, bukan hanya ujungnya}
\begin{minipage}[t]{0.494\textwidth}
\begin{icardgood}
\cardh{Uji tiap serah-terima sendiri}
Beri agen kedua serah-terima buatan dan periksa apakah ia bisa bekerja. Ini uji unit, deterministik, tanpa model.
\end{icardgood}
\end{minipage}%
\hfill
\begin{minipage}[t]{0.494\textwidth}
\begin{icardgood}
\cardh{Uji tiap agen sendiri}
Dengan alat tiruan. Kegagalan jadi bisa diatribusikan sebelum sistemnya dirangkai.
\end{icardgood}
\end{minipage}%


\vskip 4pt

\begin{minipage}[t]{0.494\textwidth}
\begin{icard}
\cardh{Uji ujung-ke-ujung, lebih jarang}
Mahal dan lambat, dan hanya berguna kalau dua uji di atas sudah hijau.
\end{icard}
\end{minipage}%
\hfill
\begin{minipage}[t]{0.494\textwidth}
\begin{icard}
\cardh{Uji kegagalan sebagian}
Apa yang terjadi kalau agen kedua gagal? Sistem harus punya jawaban, bukan menggantung.
\end{icard}
\end{minipage}%
{\fontsize{9.4}{12.6}\selectfont\color{fg2}Kartu keempat yang paling sering hilang. Pada satu agen, kegagalan berarti proses berhenti. Pada tiga agen, kegagalan di tengah bisa meninggalkan \textbf{pekerjaan setengah jadi yang sudah punya efek.}\par}\vskip 4pt
\end{frame}

\framekicker{Menguji}
\begin{frame}[shrink=25]{Kegagalan sebagian adalah masalah yang baru muncul di sini}
{\fontsize{9.4}{12.6}\selectfont\color{fg2}Satu agen yang berhenti di tengah biasanya belum melakukan apa-apa yang permanen. Tiga agen berantai bisa berhenti \textbf{sesudah} agen pertama menulis sesuatu.\par}\vskip 4pt
\begin{isteps}
\item \{'h': 'Tunda semua tulisan sampai akhir, kalau bisa', 'p': 'Kumpulkan niat, lakukan di satu titik. Ini menghapus seluruh kelas masalah.'\}
\item \{'h': 'Kalau tidak bisa, sediakan pembatalan', 'p': 'Dan ujilah pembatalannya --- jalur yang tidak pernah diuji adalah jalur yang tidak berfungsi.'\}
\item \{'h': 'Catat sampai mana yang sudah terjadi', 'p': 'Supaya manusia yang menyelesaikan tahu harus mulai dari mana.'\}
\end{isteps}
\begin{iband}[signal]
Ini alasan lain kenapa demo kredit menaruh \textbf{satu} alat tulis di ujung: pekerjaan setengah jadi yang punya efek adalah bentuk kegagalan yang paling mahal dipulihkan.
\end{iband}
\end{frame}

\framekicker{Tata kelola}
\begin{frame}[shrink=25]{Anggaran harus ada di tingkat sistem, bukan per agen}
{\fontsize{9.4}{12.6}\selectfont\color{fg2}Tiga agen dengan anggaran delapan langkah masing-masing adalah sistem dengan anggaran dua puluh empat langkah --- dan hampir tidak ada yang menuliskannya seperti itu.\par}\vskip 4pt
\begin{minipage}[t]{0.494\textwidth}
\begin{icardaccent}
\cardh{Anggaran menyeluruh}
Batas langkah, token, biaya, dan waktu untuk \textbf{seluruh proses}, dibagikan ke agen-agennya.
\end{icardaccent}
\end{minipage}%
\hfill
\begin{minipage}[t]{0.494\textwidth}
\begin{icardaccent}
\cardh{Satu tombol berhenti}
Menghentikan proses harus menghentikan semua agennya, termasuk yang sedang berjalan sejajar.
\end{icardaccent}
\end{minipage}%
\begin{iband}[signal]
Tanpa keduanya, batas yang Anda kira sudah dipasang sebenarnya \textbf{dikalikan jumlah agen} --- dan itu ditemukan pada tagihan, bukan pada rancangan.
\end{iband}
\end{frame}

\framekicker{Riset}
\begin{frame}[shrink=25]{Menyatukan temuan dari beberapa pencarian}
{\fontsize{9.4}{12.6}\selectfont\color{fg2}Bagian tersulit dari agen riset bukan mencarinya, melainkan \textbf{menyatukan} hasilnya: beberapa temuan yang sebagian tumpang tindih, sebagian bertentangan, dan sebagian dari sumber yang jauh lebih lemah daripada lainnya.\par}\vskip 4pt
\begin{isteps}
\item \{'h': 'Bawa sumbernya sampai akhir', 'p': 'Temuan tanpa asal-usul tidak bisa ditimbang, dan tidak bisa diperiksa pembaca.'\}
\item \{'h': 'Sebutkan pertentangan, jangan diratakan', 'p': 'Dua sumber yang berbeda kesimpulan adalah informasi. Ringkasan yang memilih salah satunya diam-diam membuang informasi itu.'\}
\item \{'h': 'Bedakan tidak ditemukan dari tidak ada', 'p': '``Tidak ada hasil'' sering berarti kata pencariannya yang kurang tepat --- dan itu kesimpulan yang sangat berbeda.'\}
\end{isteps}
\begin{iband}[signal]
Langkah kedua yang membedakan ringkasan riset yang berguna dari prosa yang meyakinkan: \textbf{yang berguna menyebutkan di mana sumbernya tidak sepakat.}
\end{iband}
\end{frame}

\coursesection{07}{Menjalankannya}{Menelusuri kegagalan, dan angka yang berubah artinya.}

\framekicker{Praktik}
\begin{frame}[shrink=25]{Menelusuri kegagalan pada tiga agen}
\begin{center}
\includegraphics[width=0.92\linewidth,height=0.60\textheight,keepaspectratio]{figs/agents08-debug.tex.pdf}
\end{center}
\vskip 3pt{\color{ink3}\fontsize{7.4}{9.4}\selectfont Dua pertanyaan yang menyempitkan masalah --- kalau jejaknya ada per agen.\par}
{\fontsize{9.4}{12.6}\selectfont\color{fg2}Pada satu agen, jejaknya satu berkas dan urutannya jelas. Pada tiga agen, pertanyaan pertama selalu sama: \textbf{agen mana yang terakhir punya informasi yang benar?} Kalau informasinya benar di agen 1 dan salah di agen 2, serah-terimanya yang bermasalah.\par}\vskip 4pt
\begin{iband}[signal]
Karena itu jejak per agen \textbf{plus isi tiap serah-terima} bukan kemewahan pada sistem banyak agen --- tanpa keduanya, satu-satunya cara menelusuri adalah menjalankan ulang dan menebak.
\end{iband}
\end{frame}

\framekicker{Praktik}
\begin{frame}[shrink=25]{Angka yang berubah artinya begitu ada lebih dari satu agen}
\begin{itable}{>{\raggedright\arraybackslash}p{0.2256\textwidth}>{\raggedright\arraybackslash}p{0.3384\textwidth}>{\raggedright\arraybackslash}p{0.376\textwidth}}
\thead{Angka} & \thead{Pada satu agen} & \thead{Pada banyak agen} \\ \midrule
Giliran per tugas & Panjang penalaran & Harus dipecah per agen, atau tidak berarti apa-apa \\
Biaya per tugas & Langsung & Jumlahkan semuanya, termasuk penyelia \\
Laju berhasil & Satu angka & Bisa tinggi di tiap agen dan rendah untuk sistemnya --- 0,9$^{3}$ = 0,73 \\
Waktu dinding & Jumlah langkah & Bergantung apakah benar-benar sejajar \\
\end{itable}
\begin{iband}[signal]
Baris ketiga yang paling sering mengejutkan: \textbf{tiga agen yang masing-masing berhasil 90\% menghasilkan sistem yang berhasil 73\%.} Bab 7 menghitung bentuk yang sama dengan nama pass\textasciicircum{}k.
\end{iband}
\end{frame}

\framekicker{Praktik}
\begin{frame}[shrink=25]{Urutan yang menghemat perdebatan}
\begin{isteps}
\item \{'h': 'Mulai dengan satu agen, selalu', 'p': 'Sampai ia benar-benar tidak cukup, dan sampai jejaknya menunjukkan \textbf{kenapa}.'\}
\item \{'h': 'Kalau memecah, mulai dari batas izin', 'p': 'Pemecahan yang alasannya keamanan tidak perlu dibela dengan angka --- ia menghasilkan jaminan yang tidak bisa didapat cara lain.'\}
\item \{'h': 'Ukur serah-terimanya sejak hari pertama', 'p': 'Ukurannya, dan apa yang hilang di dalamnya. Ini variabel yang menentukan biaya dan mutu sekaligus.'\}
\item \{'h': 'Bandingkan dengan satu agen yang punya semua alat', 'p': 'Di repo demo, \inlinecode{agentdemo compare research\_review} melakukan persis ini dan mencetak selisih tokennya.'\}
\end{isteps}
\begin{iband}[signal]
Langkah keempat yang membuat perdebatan selesai: \textbf{jalankan keduanya dan cetak angkanya.} Pemecahan yang benar akan menang pada angka; yang salah akan kalah, dan itu selesai dalam satu sore.
\end{iband}
\end{frame}

\framekicker{Praktik}
\begin{frame}[shrink=25]{Menggabungkan kembali jadi satu agen}
{\fontsize{9.4}{12.6}\selectfont\color{fg2}Pemecahan bukan keputusan yang tidak bisa dibatalkan, dan menggabungkan kembali sering merupakan perbaikan terbesar yang tersedia --- tetapi jarang dipertimbangkan, sebab ia terasa seperti mundur.\par}\vskip 4pt
\begin{minipage}[t]{0.494\textwidth}
\begin{icard}
\cardh{Tanda sudah waktunya digabung}
Serah-terimanya hampir sepanjang percakapannya; agen kedua sering meminta ulang; daftar alat keduanya bertumpang tindih.
\end{icard}
\end{minipage}%
\hfill
\begin{minipage}[t]{0.494\textwidth}
\begin{icardgood}
\cardh{Yang biasanya membaik setelah digabung}
Mutu naik (tidak ada yang hilang di serah-terima), penelusuran jadi sederhana, dan satu anggaran menggantikan tiga.
\end{icardgood}
\end{minipage}%
\begin{iband}[signal]
Ukurannya sama seperti waktu memutuskan memecah: \textbf{jalankan keduanya pada kumpulan uji yang sama dan cetak tiga angkanya.} Keputusan yang diambil dengan angka tidak perlu diperdebatkan dua kali.
\end{iband}
\end{frame}

\framekicker{Praktik}
\begin{frame}[shrink=25]{Lima kegagalan khas, dan tandanya}
\begin{itable}{>{\raggedright\arraybackslash}p{0.2444\textwidth}>{\raggedright\arraybackslash}p{0.3572\textwidth}>{\raggedright\arraybackslash}p{0.3384\textwidth}}
\thead{Kegagalan} & \thead{Tandanya} & \thead{Perbaikannya} \\ \midrule
Serah-terima membengkak & Ukurannya naik tiap rilis & Bentuk terstruktur dengan medan tetap \\
Penyelia jadi leher botol & Konteks penyelia paling besar & Jadikan penyelianya kode kalau urutannya tetap \\
Kerja ganda & Dua agen memanggil alat yang sama & Hasil dibagikan lewat rujukan, bukan diambil ulang \\
Tidak ada yang menghentikan & Proses panjang tanpa pemilik & Anggaran menyeluruh + satu tombol berhenti \\
Kesalahan menular & Agen 1 salah, sisanya membangun di atasnya & Pemeriksaan di serah-terima, bukan hanya di ujung \\
\end{itable}
{\fontsize{9.4}{12.6}\selectfont\color{fg2}Kelimanya terlihat pada dua hal yang sama: \textbf{jejak per agen dan isi tiap serah-terima.} Kalau keduanya ada, menelusuri sistem banyak agen tidak jauh lebih sulit daripada satu agen.\par}\vskip 4pt
\end{frame}

\framekicker{Praktik}
\begin{frame}[shrink=25]{Apa yang dilaporkan tentang sistem banyak agen}
\begin{isteps}
\item \{'h': 'Laju berhasil SISTEM, bukan per agen', 'p': 'Tiga agen 90\% menghasilkan 73\%. Angka per agen terdengar bagus dan tidak menjawab pertanyaan siapa pun.'\}
\item \{'h': 'Biaya total, termasuk penyelia', 'p': 'Koordinasi sering tidak masuk hitungan karena tidak terasa seperti pekerjaan.'\}
\item \{'h': 'Ukuran serah-terima dari waktu ke waktu', 'p': 'Ini peringatan dini yang paling murah untuk biaya dan mutu sekaligus.'\}
\item \{'h': 'Perbandingan dengan satu agen', 'p': 'Pertanyaan pertama yang akan diajukan orang yang membayar, dan pantas sudah punya jawabannya.'\}
\end{isteps}
\end{frame}

\framekicker{Praktik}
\begin{frame}[shrink=25]{Yang dipakai di demo, dan alasannya}
{\fontsize{9.4}{12.6}\selectfont\color{fg2}Repo demo punya empat kasus banyak-agen, dan masing-masing ada untuk \textbf{menunjukkan satu hal yang bisa diperiksa} --- bukan untuk memperlihatkan bahwa banyak agen itu bagus.\par}\vskip 4pt
\begin{itable}{>{\raggedright\arraybackslash}p{0.282\textwidth}>{\raggedright\arraybackslash}p{0.658\textwidth}}
\thead{Kasus} & \thead{Yang ditunjukkannya} \\ \midrule
\inlinecode{parallel\_triage} & Fan-out ketika sub-tugas benar-benar tidak saling berhubungan --- satu-satunya pemecahan yang tidak perlu dibela \\
\inlinecode{research\_review} & Pembagian rapi peneliti/penulis/kritikus, \textbf{diukur} terhadap satu agen dengan alat yang sama, token dicetak \\
\inlinecode{cross\_permission} & Pemecahan sebagai batas keamanan: agen penyelidik tidak punya alat tulis di daftarnya \\
\inlinecode{escalation} & Penyerahan ke manusia: keadaan apa yang ikut, dan bagaimana pekerjaannya kembali \\
\end{itable}
\begin{iband}[signal]
Baris kedua yang paling banyak mengubah pikiran orang di kelas, sebab hasilnya dicetak: \textbf{pembagian yang paling rapi di papan tulis kalah pada tabel biaya.}
\end{iband}
\end{frame}

\framekicker{Penutup}
\begin{frame}[shrink=25]{Yang dibawa pulang dari bab ini}
\begin{isteps}
\item \{'h': 'Tiga bentuk yang membenarkan pemecahan', 'p': 'Sejajar sungguhan, batas izin, alat yang tak berhubungan. Bagan organisasi bukan salah satunya.'\}
\item \{'h': 'Biayanya ditentukan serah-terima, bukan jumlah agen', 'p': 'Serah-terima 600 token menghemat; 3 000 token membuat mahal. Dan serah-terima cenderung membengkak.'\}
\item \{'h': 'Yang benar-benar mahal adalah informasi yang hilang', 'p': 'Ia tidak muncul di tagihan mana pun, dan ia yang membuat hasilnya memburuk.'\}
\item \{'h': '0,9 $\times$ 0,9 $\times$ 0,9 = 0,73', 'p': 'Tiga agen yang bagus bisa menghasilkan sistem yang tidak.'\}
\item \{'h': 'Bandingkan dengan satu agen, dengan angka', 'p': 'Satu sore kerja, dan perdebatannya selesai.'\}
\end{isteps}
\end{frame}
\note{Kalau ada satu latihan untuk kelas: minta mereka menyebut pemecahan yang sedang mereka rencanakan, lalu tanyakan berapa besar serah-terimanya. Yang tidak bisa menjawab biasanya belum perlu memecah.}

\end{document}
